\section{Guía de monitoreo y evaluación}
\subsection{Métricas clave}
El entrenamiento de automejora debe monitorear continuamente las siguientes métricas, para poder intervenir con rapidez cuando el win rate llegue a una meseta, la hallucination aumente o el inference compute se dispare.

\begin{table}[H]
\centering
\begin{tabular}{p{5cm}p{10cm}}
\toprule
Métrica & Descripción \\
\midrule
Win rate delta & Con Alpaca Eval 2 / DeepSeek benchmarks como referencia, observar el cambio en el win rate entre el verified response y el draft. \\
Reward variance & Si la desviación estándar del reward score en el self-rewarding aumenta de forma abrupta, indica que el modelo podría estar en overthinking o hallucinating. \\
Hallucination frequency & Usar el evaluador SelfT para etiquetar la tasa de muestras con hallucination verificada, y reducir la cross-check depth cuando se supere el umbral. \\
Inference compute & Evaluar el compute adicional que genera el verified response (número de tokens × número de cross-checks), para evitar que el costo de la evaluation exceda el presupuesto. \\
\bottomrule
\end{tabular}
\caption{Métricas centrales para monitorear el entrenamiento self-rewarding.}
\end{table}

\begin{importantbox}{Recomendación sobre la interacción entre métricas}
Cuando el win rate se estanque, revisar primero el reward variance y la hallucination frequency: si la proporción de hallucination aumenta, retroceder a las Self-Instruct seeds; si el reward variance sube, depender de la meta judge calibration para converger a tiempo.
\end{importantbox}

\subsection{Alertas automatizadas y retroalimentación}
Una vez establecidos los umbrales para las métricas correspondientes, se pueden usar automated scripts para generar alert: por ejemplo, pausar la iteración cuando el reward variance supere 0.3, o activar una revisión human-in-the-loop cuando la hallucination frequency sea > 5\% durante tres veces consecutivas. Esto puede combinarse con el flujo de Self-Feeding, enviando el alert al pipeline dashboard, de modo que el equipo de operación pueda cerrar el ciclo con rapidez dentro del \texttt{self-improvement}.

\begin{knowledgebox}{Plantilla de estrategia de alertas}
1. Definir la frecuencia: calcular las métricas inmediatamente al terminar cada iteration.\newline
2. Respuesta escalonada: el warning activa un retry prompt; el alert activa un verified response adicional.\newline
3. Archivar los registros: las muestras de hallucination junto con la trayectoria del reward se escriben en \texttt{audit.csv} para el meta judge retraining futuro.
\end{knowledgebox}

\subsection{Resumen de la sección}
Con este conjunto de métricas y alertas, el equipo puede dar seguimiento de manera transparente al estado del self-rewarding durante el entrenamiento real, y en cuanto aparezca un drift, puede recalibrar con el meta judge o repetir el self-instruct prompt.
